% =============================================================================
% METADATOS Y ESTÁNDARES DEL DOCUMENTO (PDF/UA & Tagging)
% =============================================================================
% Este archivo configura los estándares de accesibilidad y metadatos PDF.
% Se carga ANTES de \documentclass para asegurar que el motor PDF se
% inicialice correctamente.

% -----------------------------------------------------------------------------
% COMPROBACIÓN DEL MOTOR: la plantilla SOLO funciona con LuaLaTeX
% -----------------------------------------------------------------------------
% Con pdfLaTeX o XeLaTeX el documento falla más adelante con errores confusos
% (típicamente «TeX capacity exceeded, sorry [main memory size=5000000]»,
% porque esos motores tienen memoria fija y el preámbulo de esta plantilla no
% cabe en ella). Abortamos aquí con un mensaje claro.
\ifdefined\directlua\else
  \newlinechar=`\^^J
  \errmessage{^^J%
  ============================================================^^J%
   ESTA PLANTILLA REQUIERE LuaLaTeX (motor actual: no es LuaTeX)^^J%
  ============================================================^^J%
   * Overleaf: Menu -> Settings -> Compiler -> LuaLaTeX^^J%
     (Overleaf ignora la linea magica "!TeX program = lualatex")^^J%
   * Local:    make    o    lualatex -shell-escape main.tex^^J%
   * latexmk:  latexmk main.tex   (usa .latexmkrc del proyecto)^^J%
  ^^J%
   Mas informacion: docs/OVERLEAF.md^^J%
  ============================================================^^J}%
  \csname @@end\endcsname
\fi

\DocumentMetadata{
  % --- Fase de pruebas de accesibilidad (tagging) ---
  % phase-I:  Etiquetado básico automático.
  % phase-II: Etiquetado avanzado (puede dar conflictos con minted/float).
  testphase={phase-I}, 
  %
  % --- Estándar PDF ---
  % Descomenta para forzar estándares específicos. 
  % ua-2 es el estándar moderno de accesibilidad universal (sugerido).
  % a-2b es común para archivado a largo plazo.
  pdfstandard=ua-2,  
  %
  % --- Versión PDF ---
  % 1.7 es el estándar más compatible. 2.0 permite características nuevas.
  % Para PDF/UA-2 se recomienda usar PDF 2.0.
  pdfversion=2.0,
  %
  % --- Idioma principal del documento ---
  % IMPORTANTE: Este valor debe coincidir con 'idioma' en configuracion.tex
  %
  % Valores según idioma en EPSsetup:
  %   idioma = espanol    →  lang=es-ES
  %   idioma = valenciano →  lang=ca-ES
  %   idioma = ingles     →  lang=en-GB
  %
  lang=es-ES
}
